\documentclass[aspectratio=169]{beamer}
\input{header}
\coursecode{JKSSB}

\title{Data Representation, Number Systems \& Logic}
\subtitle{Lesson 50 --- From Binary Digits to Gates}
\author{JKSSB Computer Awareness}
\date{\today}

\begin{document}
\frame{\titlepage}

\begin{frame}{Why Computers Use Binary}
  \begin{itemize}
    \item Digital circuits can reliably distinguish two states, commonly
      labelled 0 and 1.
    \item A \key{number system} uses a base (radix) and a set of allowed
      digits.
    \item In base $b$, each position represents a power of $b$.
  \end{itemize}
  \begin{block}{Positional value}
    $1011_2 = 1(2^3)+0(2^2)+1(2^1)+1(2^0)=11_{10}$.
  \end{block}
\end{frame}

\begin{frame}{Four Number Systems}
  \begin{center}
    \begin{tabular}{lll}
      \toprule
      \textbf{Name} & \textbf{Base} & \textbf{Digits} \\
      \midrule
      Binary & 2 & 0--1 \\
      Octal & 8 & 0--7 \\
      Decimal & 10 & 0--9 \\
      Hexadecimal & 16 & 0--9, A--F \\
      \bottomrule
    \end{tabular}
  \end{center}
  \vspace{0.3cm}
  \muted{In hexadecimal, A=10, B=11, C=12, D=13, E=14, and F=15.}
\end{frame}

\begin{frame}{Convert Decimal to Binary}
  \begin{exampleblock}{Convert $45_{10}$ to binary}
    Divide repeatedly by 2 and read the remainders from bottom to top.
  \end{exampleblock}
  \begin{itemize}
    \item $45 \div 2 = 22$ remainder 1
    \item $22 \div 2 = 11$ remainder 0; $11 \div 2 = 5$ remainder 1
    \item $5 \div 2 = 2$ remainder 1; $2 \div 2 = 1$ remainder 0
    \item $1 \div 2 = 0$ remainder 1
  \end{itemize}
  \begin{alertblock}{Result}
    $45_{10}=101101_2$.
  \end{alertblock}
\end{frame}

\begin{frame}{Binary, Octal, and Hexadecimal}
  \begin{itemize}
    \item Group binary digits in threes to convert to octal:
      $1111111_2 = 001\ 111\ 111_2 = 177_8$.
    \item Group binary digits in fours to convert to hexadecimal:
      $0111\ 1111_2 = 7F_{16}$.
    \item Verify by place value: $7F_{16}=7(16)+15=127_{10}$.
  \end{itemize}
  \muted{Pad the leftmost group with zeros when needed; do not change the value.}
\end{frame}

\begin{frame}{Text and Decimal Digits Have Encodings}
  \begin{itemize}
    \item \key{ASCII}: character-encoding standard; original ASCII uses 7-bit
      codes.
    \item \key{Unicode}: assigns code points to characters across many
      writing systems; UTF-8 is a variable-length encoding.
    \item \key{BCD}: Binary-Coded Decimal; each decimal digit is represented
      by a 4-bit code.
    \item These are not interchangeable: a character code, a Unicode
      encoding, and a decimal-digit representation solve different problems.
  \end{itemize}
\end{frame}

\begin{frame}{Representing Colour: RGB}
  \begin{itemize}
    \item RGB represents colour with red, green, and blue channel values.
    \item In a common 24-bit RGB representation, each channel has 8 bits:
      $8+8+8=24$ bits per pixel.
    \item More bits per channel can represent more levels; actual image
      formats and display systems may use other layouts.
  \end{itemize}
  \begin{block}{Keep the claim scoped}
    24-bit RGB is a common representation, not a rule for every image or
    display.
  \end{block}
\end{frame}

\begin{frame}{Basic Logic Gates}
  \begin{center}
    \begin{tabular}{cc|cccc}
      \toprule
      $A$ & $B$ & AND & OR & XOR & NAND \\
      \midrule
      0 & 0 & 0 & 0 & 0 & 1 \\
      0 & 1 & 0 & 1 & 1 & 1 \\
      1 & 0 & 0 & 1 & 1 & 1 \\
      1 & 1 & 1 & 1 & 0 & 0 \\
      \bottomrule
    \end{tabular}
  \end{center}
  \vspace{0.25cm}
  \begin{itemize}
    \item AND is 1 only when both inputs are 1.
    \item OR is 1 when at least one input is 1.
    \item XOR is 1 when the inputs differ; NAND is the negation of AND.
  \end{itemize}
\end{frame}

\begin{frame}{NOT Gate and Exam Traps}
  \begin{itemize}
    \item NOT has one input and reverses it: NOT 0 = 1; NOT 1 = 0.
    \item Do not call BCD ordinary binary for the whole number: BCD encodes
      each decimal digit separately.
    \item Do not call UTF-8 fixed-length: it is variable-length.
    \item Do not confuse RGB bits per pixel with the number of colour
      channels.
  \end{itemize}
\end{frame}

\begin{frame}{Lesson 50: Recall}
  \begin{itemize}
    \item Binary, octal, decimal, and hexadecimal use bases 2, 8, 10, and 16.
    \item Place values are powers of the base.
    \item ASCII, Unicode/UTF-8, BCD, and RGB represent different kinds of
      information.
    \item Truth tables give the output for every input combination.
  \end{itemize}
\end{frame}
\end{document}
